%============================================================
% CH1 · P2 · PROBLEM 1 — MOVE ZEROES (LeetCode 283)
%============================================================
\subsection{Problem 1.2.1: Move Zeroes}
\problemlabel{1.2.1}{ssec:ch1-move-zeroes}

\problemcard{1.2.1}{Move Zeroes}{LeetCode 283}{https://leetcode.com/problems/move-zeroes/}{\exeasy\ \lc{Easy}}{Read/write pointers: compact the non-zeros, then fill the tail}

\subsubsection*{1. Problem statement}

Given an integer array \texttt{nums}, move all $0$'s to the end of it while \textbf{keeping the relative order of the non-zero elements}. The operation must be done \textbf{in place}, without making a copy of the array; the function returns nothing and the judge inspects \texttt{nums} afterwards.

\textbf{Constraints:} $1 \le n \le 10^4$, \; $-2^{31} \le \texttt{nums}[i] \le 2^{31} - 1$. \textbf{Follow-up:} could you minimise the total number of operations done?

\begin{table}[H]
\centering\small
\begin{tabular}{@{}lll@{}}
\toprule
& \textbf{nums (input)} & \textbf{nums (after the call)} \\
\midrule
Example 1 & $[0, 1, 0, 3, 12]$ & $[1, 3, 12, 0, 0]$ \\
Example 2 & $[0]$ & $[0]$ \\
\bottomrule
\end{tabular}
\end{table}

\begin{diagrambox}
\begin{tikzpicture}
    \dparray{s}{0.75}{before}{0/aout, 1/acell, 0/aout, 3/acell, 12/acell}
    \dparray{t}{-0.45}{after}{1/akeep, 3/akeep, 12/akeep, 0/aout, 0/aout}
    \draw[dep] (s1.south) -- (t0.north);
    \draw[dep] (s3.south) -- (t1.north);
    \draw[dep] (s4.south) -- (t2.north);
    \node[note, anchor=west, text width=5.2cm, align=left] at (5.0,0.15) {the non-zeros $1, 3, 12$ slide left \emph{in the same order}; the two zeros end up in the last two cells};
\end{tikzpicture}
\figcap{Example 1. Arrows show where each non-zero element goes.}
\end{diagrambox}

\subsubsection*{2. The problem model}
\begin{modelbox}[Problem Model --- Move Zeroes]
\renewcommand{\arraystretch}{1.35}
\begin{tabularx}{\linewidth}{@{}>{\bfseries}l X@{}}
Input & one array of $n$ integers (any \texttt{int}, zeros anywhere) \\
Computed & an \emph{array}, written over the input: the non-zero subsequence followed by $z$ zeros, where $z$ is the number of zeros \\
Over & elements; a \textbf{stable} rearrangement (order of non-zeros preserved) \\
Sizes & $n \le 10^4$: $\bigO{n^2} = 10^8$ would still pass, but the intended solution is $\bigO{n}$ with $\bigO{1}$ extra memory \\
Rule & no second array --- the read/write pattern exactly \\
\end{tabularx}
\end{modelbox}

\begin{diagrambox}
\begin{tikzpicture}
    \dparray{m}{0}{nums}{0/aout, 1/akeep, 0/aout, 3/akeep, 12/akeep}
    \draw[flow] (4.7,0) -- node[lbl, above]{model} (6.1,0);
    \begin{scope}[xshift=6.4cm]
        \node[cell, fill=basefill] (k0) at (0.45,0) {1}; \node[cell, fill=basefill] (k1) at (1.35,0) {3}; \node[cell, fill=basefill] (k2) at (2.25,0) {12};
        \node[cell, fill=lightgray, text=midgray] (k3) at (3.15,0) {0}; \node[cell, fill=lightgray, text=midgray] (k4) at (4.05,0) {0};
        \draw[thick, successgreen, decorate, decoration={brace, mirror, amplitude=4pt}] ([yshift=-1.5mm]k0.south west) -- ([yshift=-1.5mm]k2.south east) node[midway, below=4pt, font=\scriptsize\bfseries, text=successgreen] {non-zero subsequence};
        \draw[thick, midgray, decorate, decoration={brace, mirror, amplitude=4pt}] ([yshift=-1.5mm]k3.south west) -- ([yshift=-1.5mm]k4.south east) node[midway, below=4pt, font=\scriptsize\bfseries, text=midgray] {$z$ zeros};
    \end{scope}
\end{tikzpicture}
\figcap{The model: keep the subsequence of non-zeros (Pattern 2 with the rule $x \ne 0$), then overwrite the tail with zeros.}
\end{diagrambox}

\subsubsection*{3. Key observations}

\begin{itemize}[leftmargin=*]
    \item \textbf{The zeros carry no information.} All zeros are equal, so we never need to \emph{move} them: it is enough to write the non-zeros to the front and then write zeros into the remaining cells. This turns ``move'' into ``remove and refill''.
    \item \textbf{Read/write pointers keep the order.} Copying each non-zero to the next free position \texttt{w} preserves their relative order, and $w \le r$ means nothing unread is overwritten (invariant of Pattern~2).
    \item \textbf{Swapping from both ends is wrong.} The partition idea ``swap a zero on the left with a non-zero on the right'' does put the zeros last, but on Example 1 it produces $[12, 1, 3, 0, 0]$: the order of the non-zeros is lost.
    \item \textbf{Erasing inside an index loop is wrong (and slow).} ``If \texttt{nums[i]} is 0, erase it and \texttt{push\_back(0)}'' skips the element that slides into position $i$: on $[0, 0, 1]$ it returns $[0, 1, 0]$. Each erase also costs $\bigO{n}$.
\end{itemize}

\begin{diagrambox}
\begin{tikzpicture}
    \begin{scope}
        \dparray{w}{0}{two ends}{12/aread, 1/acell, 3/aread, 0/aout, 0/aout}
        \node[note, anchor=north, text width=4.4cm] at (2.25,-0.5) {swap $a_0 \leftrightarrow a_4$, then $a_2 \leftrightarrow a_3$: zeros last, but $12$ now precedes $1$ and $3$ --- \wa{}};
    \end{scope}
    \begin{scope}[xshift=7.4cm]
        \dparray{x}{0}{erase}{0/aread, 1/acell, 0/aout}
        \node[note, anchor=north, text width=3.6cm] at (1.35,-0.5) {input $[0, 0, 1]$: the second zero slides into index 0 and is skipped --- \wa{}};
    \end{scope}
\end{tikzpicture}
\figcap{Outputs of the two tempting wrong solutions (both computed by running them).}
\end{diagrambox}

\subsubsection*{4. Approaches}

\begin{bruteforcebox}
\textbf{Idea.} Imitate bubble sort with the key ``is zero'': repeatedly swap a zero with a non-zero right after it. Every swap moves a zero one step right past a non-zero, so the non-zeros never change their relative order; after $n$ passes all zeros are at the end.
\begin{lstlisting}[style=cpp]
#include <bits/stdc++.h>
using namespace std;

class Solution {
public:
    void moveZeroes(vector<int>& nums) {
        int n = nums.size();
        for (int pass = 0; pass < n; pass++)
            for (int i = 0; i + 1 < n; i++)
                if (nums[i] == 0 && nums[i + 1] != 0)
                    swap(nums[i], nums[i + 1]);       // zero moves one step right
    }
};

int main() {
    vector<vector<int>> tests = {{0, 1, 0, 3, 12}, {0}};
    for (auto& t : tests) {
        Solution().moveZeroes(t);
        for (int x : t) cout << x << " ";            // 1 3 12 0 0 | 0
        cout << "| ";
    }
    cout << "\n";
}
\end{lstlisting}
\begin{complexitybox}
\textbf{Time:} $\bigO{n^2}$ --- $10^8$ comparisons for $n = 10^4$: borderline \ac{} on LeetCode, far too slow for larger inputs. \textbf{Space:} $\bigO{1}$. The obvious alternative --- copy the non-zeros into a new vector --- is $\bigO{n}$ but breaks the ``no copy'' rule.
\end{complexitybox}
\end{bruteforcebox}

\begin{algorithmbox}
\textbf{Idea.} Pattern 2 with the keep rule $x \ne 0$: a write pointer \texttt{w} starts at 0; the read pointer \texttt{r} scans the array and copies every non-zero to \texttt{nums[w++]}. Afterwards \texttt{nums[0..w)} holds the non-zeros in order, and the cells \texttt{nums[w..n)} --- which contain leftovers --- are overwritten with zeros.
\begin{lstlisting}[style=cpp]
#include <bits/stdc++.h>
using namespace std;

class Solution {
public:
    void moveZeroes(vector<int>& nums) {
        int n = nums.size(), w = 0;
        for (int r = 0; r < n; r++)
            if (nums[r] != 0) nums[w++] = nums[r];    // compact the non-zeros (stable)
        while (w < n) nums[w++] = 0;                  // the tail becomes zeros
    }
};

int main() {
    vector<vector<int>> tests = {{0, 1, 0, 3, 12}, {0}, {0, 0, 1}, {1, 0, 2, 0, 0, 3}};
    for (auto& t : tests) {
        Solution().moveZeroes(t);
        for (int x : t) cout << x << " ";   // 1 3 12 0 0 | 0 | 1 0 0 | 1 2 3 0 0 0
        cout << "| ";
    }
    cout << "\n";
}
\end{lstlisting}
\begin{complexitybox}
\textbf{Time:} $\bigO{n}$ --- one pass of $n$ reads plus exactly $n$ writes in total ($w$ copies and $n - w$ zero fills). \textbf{Space:} $\bigO{1}$. \textbf{Verdict:} \ac{}.
\end{complexitybox}
\end{algorithmbox}

\begin{invariantbox}[Why it is correct]
\textbf{Loop invariant} (Pattern 2): before the step with index $r$, $w \le r$, the cells \texttt{nums[0..w)} hold the non-zero elements of the original \texttt{nums[0..r)} in their original order, and \texttt{nums[r..n)} is untouched. The proof is the general one of the pattern theory: a non-zero is written to position $w \le r$, i.e.\ to a cell already read, and appended to the kept list.

\textbf{After the loop} ($r = n$), \texttt{nums[0..w)} is the full non-zero subsequence, so $w$ equals the number of non-zeros and $n - w = z$, the number of zeros. \textbf{The fill} writes $z$ zeros into \texttt{nums[w..n)}. The final array is ``non-zeros in order, then $z$ zeros'', which is exactly the required arrangement. No second array was used. \hfill$\blacksquare$
\end{invariantbox}

\begin{alternativebox}
\textbf{Swap instead of copy --- one pass, no fill.} When \texttt{nums[r]} is non-zero, \emph{swap} it with \texttt{nums[w]} instead of overwriting. Then the cells in $[w, r)$ are always zeros (each swap sends a zero to position $r$), and when the loop ends the tail is already zero.
\begin{lstlisting}[style=cpp]
#include <bits/stdc++.h>
using namespace std;

class Solution {
public:
    void moveZeroes(vector<int>& nums) {
        int w = 0;                                   // invariant: nums[w..r) are all zeros
        for (int r = 0; r < (int)nums.size(); r++)
            if (nums[r] != 0) {
                if (r != w) swap(nums[w], nums[r]);  // skip useless self-swaps
                w++;
            }
    }
};

int main() {
    vector<vector<int>> tests = {{0, 1, 0, 3, 12}, {0}, {0, 0, 1}, {1, 0, 2, 0, 0, 3}};
    for (auto& t : tests) {
        Solution().moveZeroes(t);
        for (int x : t) cout << x << " ";   // 1 3 12 0 0 | 0 | 1 0 0 | 1 2 3 0 0 0
        cout << "| ";
    }
    cout << "\n";
}
\end{lstlisting}
\textbf{Why $[w, r)$ holds only zeros:} it is empty at the start; a zero at $r$ extends it by one cell; a non-zero at $r$ is swapped with the first cell of the block (a zero, or itself when the block is empty), which moves the block one step right. It is the same algorithm in disguise, with a different write count: each useful swap costs two writes, but nothing is written when the array has no zeros before its non-zeros. This answers the follow-up question:

\begin{table}[H]
\centering\small
\begin{tabular}{@{}lcc@{}}
\toprule
\textbf{array writes} & \textbf{copy + fill} & \textbf{swap} \\
\midrule
$[0, 1, 0, 3, 12]$ (Example 1) & 5 & 6 \\
$[1, 2, 3, 0]$ (zeros already last) & 4 & 0 \\
general & always $n$ & $2 \times$ (non-zeros that are not yet in place) \\
\bottomrule
\end{tabular}
\end{table}

\begin{diagrambox}
\begin{tikzpicture}
    \dparray{q}{0}{\texttt{nums}}{1/akeep, 3/akeep, 0/aout, 0/aout, 12/aread}
    \aptr{q2}{w}{accentorange}
    \aptr{q4}{r}{successgreen}
    \draw[dep] (q4.north) to[bend right=35] node[lbl, above]{swap} (q2.north);
    \node[note, anchor=west, text width=4.4cm, align=left] at (5.0,0) {Example 1 at $r = 4$: the zero block $[2, 4)$ moves one step right after the swap; final array $[1, 3, 12, 0, 0]$};
\end{tikzpicture}
\figcap{The swap version keeps a block of zeros between \texttt{w} and \texttt{r} and rolls it to the end.}
\end{diagrambox}
\end{alternativebox}

\begin{tricksbox}
\begin{itemize}[leftmargin=*]
    \item \textbf{Library one-liner:} \texttt{fill(stable\_partition(...))} works but allocates a buffer; \texttt{auto e = remove(nums.begin(), nums.end(), 0); fill(e, nums.end(), 0);} is exactly Approach 2 (\texttt{std::remove} is a stable read/write compaction).
    \item \textbf{Fewer writes in the copy version:} write only when $r \ne w$ (\texttt{if (r != w) nums[w] = nums[r];}). Then an array without zeros causes no writes at all, and the fill writes exactly $z$ cells.
    \item \textbf{Counting is enough here} because all removed values are identical. If the removed elements were different (say ``move all negative numbers to the end, keeping both groups in order''), the fill would not be possible, and an $\bigO{1}$-memory stable solution becomes much harder --- a reason to accept $\bigO{n}$ extra memory there.
\end{itemize}
\end{tricksbox}

\subsubsection*{5. Trace}

\begin{tracebox}[{Example 1: $\texttt{nums} = [0, 1, 0, 3, 12]$, copy + fill}]
\centering
\begin{tikzpicture}
    \begin{scope}
        \dparray{a}{0}{$r = 1$}{1/awrite, 1/aread, 0/acell, 3/acell, 12/acell}
        \aptr{a1}{r}{successgreen}
        \node[note, anchor=west, text width=4.0cm, align=left] at (4.8,0) {$a_1 = 1 \ne 0$: copied to $a_0$, $w = 1$ ($r = 0$ was a zero: skipped)};
    \end{scope}
    \begin{scope}[yshift=-1.75cm]
        \dparray{b}{0}{$r = 3$}{1/akeep, 3/awrite, 0/aout, 3/aread, 12/acell}
        \aptr{b3}{r}{successgreen}
        \node[note, anchor=west, text width=4.0cm, align=left] at (4.8,0) {$a_2 = 0$ skipped; $a_3 = 3$ copied to $a_1$, $w = 2$};
    \end{scope}
    \begin{scope}[yshift=-3.5cm]
        \dparray{c}{0}{$r = 4$}{1/akeep, 3/akeep, 12/awrite, 3/aout, 12/aread}
        \aptr{c4}{r}{successgreen}
        \node[note, anchor=west, text width=4.0cm, align=left] at (4.8,0) {$12$ copied to $a_2$, $w = 3$; cells $3, 4$ hold leftovers};
    \end{scope}
    \begin{scope}[yshift=-5.25cm]
        \dparray{d}{0}{fill}{1/adone, 3/adone, 12/adone, 0/awrite, 0/awrite}
        \node[note, anchor=west, text width=4.0cm, align=left] at (4.8,0) {\texttt{while (w < n)} writes the two zeros};
    \end{scope}
\end{tikzpicture}
\figcap{Blue: the cell written in this step. Orange: the cell being read. The leftovers $3, 12$ in cells $3$ and $4$ are harmless: they were already read.}

\medskip
\small
\begin{tabular}{@{}cccll@{}}
\toprule
$r$ & \texttt{nums[r]} & \textbf{action} & $w$ after & \textbf{array after the step} \\
\midrule
0 & 0 & skip & 0 & $[0, 1, 0, 3, 12]$ \\
1 & 1 & \texttt{nums[0] = 1} & 1 & $[1, 1, 0, 3, 12]$ \\
2 & 0 & skip & 1 & $[1, 1, 0, 3, 12]$ \\
3 & 3 & \texttt{nums[1] = 3} & 2 & $[1, 3, 0, 3, 12]$ \\
4 & 12 & \texttt{nums[2] = 12} & 3 & $[1, 3, 12, 3, 12]$ \\
fill & & \texttt{nums[3..5) = 0} & 5 & $[1, 3, 12, 0, 0]$ \\
\bottomrule
\end{tabular}
\end{tracebox}

\begin{tracebox}[{Copy + fill on $[1, 0, 2, 0, 0, 3]$}]
\centering\small
\begin{tabular}{@{}cccll@{}}
\toprule
$r$ & \texttt{nums[r]} & \textbf{action} & $w$ after & \textbf{array after the step} \\
\midrule
0 & 1 & \texttt{nums[0] = 1} (self-copy) & 1 & $[1, 0, 2, 0, 0, 3]$ \\
1 & 0 & skip & 1 & $[1, 0, 2, 0, 0, 3]$ \\
2 & 2 & \texttt{nums[1] = 2} & 2 & $[1, 2, 2, 0, 0, 3]$ \\
3, 4 & 0, 0 & skip & 2 & $[1, 2, 2, 0, 0, 3]$ \\
5 & 3 & \texttt{nums[2] = 3} & 3 & $[1, 2, 3, 0, 0, 3]$ \\
fill & & \texttt{nums[3..6) = 0} & 6 & $[1, 2, 3, 0, 0, 0]$ \\
\bottomrule
\end{tabular}
\end{tracebox}

\subsubsection*{6. Test suite}

\begin{testbox}
\centering\small
\begin{tabular}{@{}lll@{}}
\toprule
\textbf{nums} & \textbf{Why it is interesting} & \textbf{Expected} \\
\midrule
$[0]$ & official Example 2: a single zero & $[0]$ \\
$[1]$ & a single non-zero & $[1]$ \\
$[0, 0, 1]$ & consecutive zeros first (erase-loop trap) & $[1, 0, 0]$ \\
$[1, 2, 3]$ & no zeros: nothing moves & $[1, 2, 3]$ \\
$[0, 0, 0]$ & only zeros & $[0, 0, 0]$ \\
$[4, 0, 0, -2^{31}, 5, 0]$ & negative extreme value, zero at the end & $[4, -2^{31}, 5, 0, 0, 0]$ \\
$[1, 0, 2, 0, 0, 3]$ & interleaved zeros & $[1, 2, 3, 0, 0, 0]$ \\
$[0, 2, 0, 4, \dots, 0, 10000]$ & $n = 10^4$, every other element zero & $[2, 4, \dots, 10000, 0, \dots, 0]$ \\
\bottomrule
\end{tabular}
\end{testbox}
Both linear solutions and the bubble brute force reproduce the official examples and agree with a copy-based reference on 600 random arrays ($n \le 9$, values in $[-3, 3]$, so about one element in seven is a zero).

\subsubsection*{7. Common mistakes}

\begin{warningbox}[Mistake 1 --- forgetting the fill]
Without \texttt{while (w < n) nums[w++] = 0;} Example 1 ends as $[1, 3, 12, 3, 12]$: the leftovers stay. The compaction alone is the answer only for problems that return a new length (Remove Element, Problem~\ref{ssec:ch1-remove-duplicates-ii}).
\end{warningbox}

\begin{warningbox}[Mistake 2 --- partitioning from both ends]
Fast and in place, but not stable: Example 1 becomes $[12, 1, 3, 0, 0]$. Read the words ``relative order''.
\end{warningbox}

\begin{warningbox}[Mistake 3 --- \texttt{erase} in an index loop]
\begin{lstlisting}[style=snippet]
for (int i = 0; i < n; i++)
    if (nums[i] == 0) { nums.erase(nums.begin() + i); nums.push_back(0); }   // WRONG
\end{lstlisting}
$[0, 0, 1] \to [0, 1, 0]$: after the erase, the next element sits at index $i$ and is skipped. Even with \texttt{i--} it is $\bigO{n^2}$.
\end{warningbox}

\begin{warningbox}[Mistake 4 --- a copy ``to be safe'']
\texttt{vector<int> tmp; \ldots{} nums = tmp;} passes LeetCode's checker but violates the statement. In an interview it is the answer to a different question.
\end{warningbox}

\subsubsection*{8. Takeaways}

\begin{itemize}[leftmargin=*]
    \item ``Move all $X$ to the end, keep the order of the rest'' $=$ compact the rest with read/write pointers, then fill the tail with $X$ (possible because all $X$ are equal).
    \item The swap version keeps a block of $X$'s between \texttt{w} and \texttt{r}; it needs no fill and writes nothing when no element has to move.
    \item Stability decides the technique: read/write compaction is stable, swapping from both ends is not.
    \item $\bigO{1}$ extra memory and $\bigO{n}$ time are both achievable; \texttt{erase} in a loop achieves neither.
\end{itemize}

\begin{variantsbox}
\variant{Remove a value, return the length}{\exeasy\ \lc{Easy}\ \varlink{Remove Element}{LeetCode 27}{https://leetcode.com/problems/remove-element/}}
  {remove every occurrence of \texttt{val} in place and return the number $k$ of remaining elements; order may change.}
  {the subsequence of elements $\ne \mathit{val}$; only \texttt{nums[0..k)} is judged.}
  {the compaction loop without the fill, return $w$. $[3,2,2,3]$, val $3 \to k = 2$ $[2,2]$; $[0,1,2,2,3,0,4,2]$, val $2 \to k = 5$.}
\variant{Evens before odds}{\exeasy\ \lc{Easy}\ \varlink{Sort Array By Parity}{LeetCode 905}{https://leetcode.com/problems/sort-array-by-parity/}}
  {put all even numbers before all odd numbers; any order inside the groups is accepted.}
  {two groups, no stability required --- the removed elements (odds) are \emph{different}, so a fill is impossible.}
  {the swap version with ``keep $=$ even'' moves the odds to the back; $[3,1,2,4] \to [2,4,3,1]$.}
\variant{Merge equal neighbours, then move zeroes}{\exeasy\ \lc{Easy}\ \varlink{Apply Operations to an Array}{LeetCode 2460}{https://leetcode.com/problems/apply-operations-to-an-array/}}
  {first, for $i = 0..n-2$ in order: if $a_i = a_{i+1}$ then $a_i \gets 2a_i$, $a_{i+1} \gets 0$; then move the zeroes.}
  {a left-to-right simulation followed by this problem.}
  {one simulation pass, then copy + fill; $[1,2,2,1,1,0] \to [1,4,2,0,0,0]$, $[0,1] \to [1,0]$.}
\variant{Zeros to the front}{\varnone}
  {move all zeros to the \emph{front}, keeping the order of the non-zeros.}
  {the mirror image: the non-zeros are compacted towards the right end.}
  {run $r$ and $w$ from $n-1$ down to $0$, then fill $a[0..w]$ with zeros; $[0,1,0,3,12] \to [0,0,1,3,12]$.}
\end{variantsbox}
